\documentclass[12pt,a4paper]{article}
\usepackage[utf8]{inputenc}
\usepackage{amsmath,amssymb,amsthm}
\usepackage{algorithm}
\usepackage{algorithmic}
\usepackage{geometry}
\geometry{margin=1in}

\newtheorem{theorem}{Theorem}
\newtheorem{lemma}[theorem]{Lemma}

\title{Problem 63: Powerful Digit Counts}
\author{Project Euler Solution}
\date{}

\begin{document}
\maketitle

\section{Problem Statement}

How many $n$-digit positive integers exist which are also an $n$-th power?

\section{Mathematical Foundation}

\begin{theorem}[Complete characterization]
A positive integer $a^n$ has exactly $n$ digits if and only if $a \in \{1, 2, \ldots, 9\}$ and $1 \le n \le N(a)$, where
\[
N(a) = \begin{cases} 1 & \text{if } a = 1, \\ \left\lfloor \dfrac{1}{1 - \log_{10} a} \right\rfloor & \text{if } 2 \le a \le 9. \end{cases}
\]
No solution exists for $a \ge 10$ or $a = 0$.
\end{theorem}

\begin{proof}
The condition for $a^n$ to have exactly $n$ digits is
\[
10^{n-1} \le a^n < 10^n. \tag{$\star$}
\]

\textbf{Step 1: Bounding $a$.} The right inequality gives $a < 10$, so $a \in \{1, \ldots, 9\}$. For $a \ge 10$: $a^n \ge 10^n$, yielding $\ge n+1$ digits.

\textbf{Step 2: Bounding $n$.} Taking $\log_{10}$ of the left inequality:
\[
n - 1 \le n \log_{10} a \implies n(1 - \log_{10} a) \le 1.
\]

For $a = 1$: the constraint becomes $n \le 1$, and only $1^1 = 1$ works.

For $2 \le a \le 9$: since $0 < 1 - \log_{10} a < 1$, we get $n \le 1/(1 - \log_{10} a)$, so
\[
N(a) = \left\lfloor \frac{1}{1 - \log_{10} a} \right\rfloor.
\]
The right inequality $a^n < 10^n$ is automatic since $a < 10$. The left inequality holds for all $n \le N(a)$ because $n(1 - \log_{10} a) \le 1$ implies $10^{n-1} \le a^n$.
\end{proof}

\begin{lemma}[Explicit enumeration]
The values $N(a)$ for $a = 1, \ldots, 9$ are:
\begin{center}
\begin{tabular}{cccc}
$a$ & $\log_{10} a$ & $1/(1 - \log_{10} a)$ & $N(a)$ \\
\hline
1 & 0.0000 & 1.000 & 1 \\
2 & 0.3010 & 1.431 & 1 \\
3 & 0.4771 & 1.912 & 1 \\
4 & 0.6021 & 2.513 & 2 \\
5 & 0.6990 & 3.322 & 3 \\
6 & 0.7782 & 4.507 & 4 \\
7 & 0.8451 & 6.456 & 6 \\
8 & 0.9031 & 10.319 & 10 \\
9 & 0.9542 & 21.854 & 21 \\
\end{tabular}
\end{center}
\end{lemma}

\begin{proof}
Direct computation. Each entry is verified: $a^{N(a)}$ has $N(a)$ digits, $a^{N(a)+1}$ does not have $N(a)+1$ digits.
\end{proof}

\begin{theorem}[Total count]
The number of $n$-digit $n$-th powers is
\[
\sum_{a=1}^{9} N(a) = 1 + 1 + 1 + 2 + 3 + 4 + 6 + 10 + 21 = 49.
\]
\end{theorem}

\begin{proof}
Sum the values from the lemma. By Theorem~1, this enumeration is exhaustive over all positive integers and all exponents.
\end{proof}

\section{Algorithm}

\begin{algorithm}[H]
\caption{Count $n$-digit $n$-th powers}
\begin{algorithmic}[1]
\STATE $\mathrm{count} \gets 0$
\FOR{$a = 1$ \TO $9$}
    \IF{$a = 1$}
        \STATE $N \gets 1$
    \ELSE
        \STATE $N \gets \lfloor 1 / (1 - \log_{10} a) \rfloor$
    \ENDIF
    \STATE $\mathrm{count} \gets \mathrm{count} + N$
\ENDFOR
\RETURN count
\end{algorithmic}
\end{algorithm}

\section{Complexity Analysis}

\textbf{Time:} $O(1)$. Exactly 9 iterations with constant-time arithmetic each.

\textbf{Space:} $O(1)$. Only scalar variables.

\section{Answer}

\[
\boxed{49}
\]

\end{document}
